\section{Metodología de la Investigación}

Para el desarrollo del trabajo de grado, se propuso un plan de 3 fases, que acarrea una secuencia lógica desde la implementación hasta el análisis de los resultados:


\subsection{Implementación de RK4}

Esta primera fase se basó en el análisis de la fuente de código inicial provista por el director, que contiene una implementación del algoritmo RK4 para la visualización de Streamlines, Streaklines y Pathlines usando CUDA. Posterior a ello, la base de código fue modificada para facilitar su estudio y migración a demás plataformas que usan SYCL o HIP para su estudio en fases próximas.

\subsubsection*{Actividades}

\begin{enumerate}
	\itemsep-2mm
	\item Recopilación de datos de entrada de tamaño variable para pruebas.
	\item Implementación del algoritmo RK4 en CUDA.
	\item Implementación del algoritmo RK4 en SYCL.
\end{enumerate}

\subsubsection*{Productos}

\begin{enumerate}
	\itemsep-2mm
	\item Implementaciones del algoritmo RK4 para visualización de Streamlines, Streaklines y Pathlines en los siguientes estándares: CUDA, HIP, oneAPI (SYCL*\footnote{Denominación que otorga Intel a su uso de SYCL, pues es una extensión del mismo}) y SYCL
	\item Evaluación de los posibles lenguajes de programación para la implementación.
\end{enumerate}

\subsection{Desarrollo del ambiente de pruebas}

La segunda fase consistió en la definición del cómo se realiza la declaración de la arquitectura. Partiendo de los criterios de selección establecidos en la fase 1, se determinó un lenguaje de notación el cual permitió definir la arquitectura objetivo a alcanzar, al igual que la gramática correspondiente para poder realizar dicha declaración.

La segunda fase consistió de la creación de ambientes controlados donde se ejecutarán los algoritmos, manteniendo las mismas condiciones de cómputo para los algoritmos, y se recopilarán las métricas establecidas para la comparación. Donde se usó Docker y la misma imágen base de Ubuntu para instalar las dependencias de cada uno de los estándares a evaluar. Para esto se hizo usó del clúster de SC3 para ejecutar las imágenes.

\subsubsection*{Actividades}

\begin{enumerate}
	\itemsep-2mm
	\item Determinación de métricas a evaluar.
	\item Implementación en Docker del ambiente de pruebas.
	\item Ejecución síncrona de pruebas para ambas implementaciones.
\end{enumerate}

\subsubsection*{Productos}

\begin{enumerate}
	\itemsep-2mm
	\item Listado de métricas que se recopilarán y almacenarán durante la ejecución de las imágenes de Docker
	\item Dockerfiles correspondientes a cada uno de los 4 estándares a evaluar
	\item Archivos csv con los datos recopilados de las pruebas
\end{enumerate}

\subsection{Recopilación de resultados}

Durante la tercera fase del proyecto, se buscó determinar e implementar cómo se realiza la comparación entre el estado de la arquitectura obtenido durante la auto-descripción de la misma y el objetivo establecido. Así mismo, y con el fin de reportar a los administradores de los sistemas, también fue necesario definir los \textit{niveles} de similitud entre las 2 arquitecturas.

Durante la tercera fase y última fase, se compilarán los resultados en un formato adecuado para su análisis, y se retroalimentarán los resultados obtendidos para concluir el trabajo de grado.

Durante la tercera y última fase, se compilaron y organizaron los resultados en formatos csv usando Python,se realizaron gráficas que muestran las diferencias entre los 4 estándares en las métricas obtenidas. De la misma forma, se incluyó el aspecto cualitativo de la dificultad de su uso, dadas las condiciones del hardware y software durante el desarrollo de la investigación como punto de comparación entre los estándares.

\subsubsection*{Actividades}

\begin{enumerate}
	\itemsep-2mm
	\item Organización de los resultados obtenidos en gráficos y tablas.
	\item Retroalimentación y conclusiones de los resultados obtenidos.
\end{enumerate}

\subsubsection*{Productos}

\begin{enumerate}
	\itemsep-2mm
	\item Tablas y gráficas correspondientes a las métricas obtenidas y su comparación para cada estándar.
	\item Evaluación de los estándares de acuerdo a los datos cuantitativos y cualitativos obtenidos.
\end{enumerate}

